\chapter{结果分析与讨论}

\section{完成工作强度与形成能力}

\figfull[.98\textwidth]{engineering_workload_dashboard.pdf}{第一年度可核验工作量总览。七组数字分别对应数据资产、正式训练数据、训练探索、正式训练、部署模型、视觉审查和强化学习研究；由于单位不同，它们用于展示工作覆盖面，不能直接相加。}

这些数字串起的是一条连续的工程链，而不是彼此孤立的成果：专用采集软件把双臂示教稳定保存下来，数据编辑中心把轨迹组织成可检查、可发布的版本，训练平台支持多轮资源与配置探索，正式保存模型进入现场运行，注意力记录帮助复核三路视野的使用情况，强化学习研究又在独立数据上形成新的训练与验证闭环。由此形成的资产可以在下一年度继续复用，工作价值不依赖某一次展示是否完美。

\section{已形成的核心能力}

\subsection{完整工程闭环已经形成}

第一年度最有分量的成果不是一条孤立的损失曲线，而是数据生产、数据管理、版本发布、大规模训练、保存模型和现场运行已经串成闭环。数据中心有效项目汇总与逐条轨迹统计一致，正式模型数据可追溯到25个固定版本，部署保存模型也能与训练历史对应。\ev{E02--E09} 这意味着后续新增瓶型、补采困难条件或比较新模型时，可以复用现有生产线，而不需要重新从零搭建。

\subsection{数据覆盖从固定动作扩展到多条件长流程}

真实示教包含单瓶、多瓶、方向变化、无盖、含水、回位、翻转和自由旋转瓶盖等条件。团队还通过重复采样提高自由旋转瓶盖场景的出现机会。这项工作把训练材料从单一标准动作扩展为覆盖多类现场变化的数据资产，也为定位空拧和倒抓提供了可追溯的数据入口。

\subsection{模型与现场系统已经完成集成}

正式模型共同使用顶部总览、左右腕部近景和双臂状态，连续生成未来50个控制时刻的动作，并在当前动作执行期间准备下一段动作。现场观察到完整取瓶、协同旋盖和分类投放循环，与这一运行链路相互印证；8,223个真实注意力样本又证明三路视野确实进入了运行审查流程。\ev{E06,E07,E10,E12}

\subsection{研究工作已经从基础模仿扩展到能力优化}

在基础模型之外，团队还完成了835条强化学习研究轨迹、238,816个训练片段、连续5个轮次和固定离线验证。它尚未取代基础模型，但已经建立了“真实数据—人工评价—动作优化—价值判断—保存与验证”的研究路径，为下一年度在统一真机协议下比较增益做好准备。

\section{任务难度、完成工作与阶段能力}

\begin{table}[H]
\centering
\caption{任务难度—完成工作—工作量—阶段能力对应表}
\begin{tabularx}{\textwidth}{>{\bfseries}p{25mm}p{39mm}p{36mm}X}
\toprule
任务难点 & 已完成工作 & 可核验工作量 & 形成的阶段能力\\
\midrule
随机瓶堆与遮挡 & 多项目、多条件、多视角采集和数据编辑 & 51个有效项目、2,413条轨迹 & 能从桌面多瓶场景中选择并接近目标\\
瓶口与方向判断 & 方向、无盖、翻转、腕部近景数据专项覆盖 & 方向404条、无盖158条、翻转127条 & 对多种摆放与瓶盖条件形成初步适应基础\\
双臂协同接触 & 完整取瓶、稳瓶、找盖、旋拧与分箱示教 & 1,051条正式训练轨迹，14维双臂状态和动作 & 已观察到完整双臂任务循环\\
长流程连续控制 & 50步动作片段、提前规划下一片段、现场集成 & 50Hz控制，正式部署第19,000步保存模型 & 已观察到超过单次动作的连续循环表现\\
大模型训练稳定性 & 多批量规模、资源和配置探索 & 41次尝试、14组配置、6,059条正式训练记录 & 建立可重复训练与保存模型链路\\
视觉使用与诊断 & 三路真实注意力记录和统计 & 9份清单、8,223个样本 & 可审查全局与腕部视野的使用情况\\
后续能力优化 & 强化学习数据、训练、保存和离线验证 & 835条轨迹、238,816个片段、5个轮次 & 形成下一阶段模型优化研究基础\\
\bottomrule
\end{tabularx}
\end{table}

这张对应表说明，工作量并非单纯堆积数据或重复训练。每一类投入都针对一个具体任务难点，并形成下一环节可继续利用的能力或工程资产。

\section{成果成熟度与下一阶段量化}

\begin{table}[H]
\centering
\caption{主要成果的成熟度与下一步量化方式}
\begin{tabularx}{\textwidth}{>{\bfseries}p{48mm}p{27mm}X}
\toprule
成果 & 当前成熟度 & 下一步量化\\
\midrule
数据、训练、保存和部署闭环 & 可复核 & 保持固定版本和自动审计\\
连续双臂动作生成与现场集成 & 可复核 & 增加端到端延迟和阶段时间记录\\
完整任务循环 & 现场已观察 & 建立逐瓶编号和阶段判定\\
超过1小时、约2瓶/分钟 & 现场观测估计 & 增加起止时间、有效工作时间和逐瓶计数\\
多条件适应与初步泛化 & 已有数据与现场迹象 & 使用冻结条件集重复测试\\
强化学习动作优化 & 离线闭环已形成 & 与基础模型开展同条件真机比较\\
\bottomrule
\end{tabularx}
\end{table}

这种分层并不削弱已经完成的工作：工程链路和完整任务能力已经形成，下一步重点是把现场表现转换成可重复、可比较、可统计的正式指标。

\section{持续提升的典型场景}

\begin{table}[H]
\centering
\caption{典型场景、已有基础与下一步强化}
\begin{tabularx}{\textwidth}{>{\bfseries}p{28mm}p{47mm}p{45mm}X}
\toprule
场景 & 已有基础 & 下一步强化 & 验收方式\\
\midrule
无盖瓶仍进入旋拧 & 已有158条无盖轨迹并定位到任务指令条件缺口 & 增加有盖/无盖配对数据和条件指令 & 记录是否正确跳过旋拧阶段\\
部分瓶子方向相反 & 已有404条方向数据和127条翻转数据 & 增加瓶口方向、抓取端和抬起结果标签 & 固定姿态集合逐次统计正确抓取端\\
接触阶段波动 & 双腕近景和完整旋拧示教已经形成 & 增加瓶盖角度与夹爪负载记录 & 分别统计抓盖、转动和脱离\\
长流程偏差积累 & 七阶段任务链和连续动作控制已经实现 & 增加阶段时间戳与异常恢复示教 & 定位每轮最早偏差并记录恢复率\\
评测结果可复核性 & 现场已经观察到长时间连续工作 & 建立固定场景、逐瓶编号和统一计时 & 输出成功率、阶段完成率和吞吐量区间\\
\bottomrule
\end{tabularx}
\end{table}

\section{两个重点问题已经形成可执行改进路径}

\subsection{无盖条件判断：从现场现象回溯到数据指令}

数据交叉检查发现，无盖画面与无条件旋盖指令同时存在。这一结果把“空拧”从笼统现象收敛到可直接操作的数据问题。下一步不需要盲目增加所有数据，而应明确增加“观察瓶口—判断有盖—决定是否旋拧”的阶段，并针对同瓶型、同位置采集有盖/无盖配对轨迹。完成后以“是否正确进入旋拧阶段”作为直接指标。

\subsection{瓶身方向：已有覆盖基础，重点转向精细标签}

方向相关轨迹数量已经达到404条，说明团队并非没有处理方向变化。当前更有价值的工作是为瓶口方向、抓取端、抬起后姿态和后续可完成性增加结构化标签，从而区分“看错方向、选错目标、夹持后滑动或纠正不及时”。这种标注能够把新增采集集中到真正薄弱的环节。

\section{训练、控制与评测体系的继续完善}

正式训练已经运行到59,990步，现场采用第19,000步保存版本。下一阶段应冻结3--5个候选保存点，在同一场景、同一次数和同一阶段判定下比较，而不是只按训练误差选择。这样既能保留现场经验，也能形成可复核的保存模型选择依据。

50Hz控制和50步动作片段为连续双臂运动提供了基础。后续记录端到端推理时间、片段切换抖动和控制丢帧，可以进一步优化异常出现后的反应速度。当前注意力导出开销约40--55毫秒中位数，只代表诊断记录的额外开销，不作为完整推理延迟。

\section{强化学习与数字孪生的研究基础}

强化学习部分已经形成真实冲洗数据、人工评价、动作优化、价值判断、保存模型和离线验证闭环；数字孪生方向正在为碰撞检查、初始姿态覆盖、插入轨迹与安全边界预演建立基础。这两条路线与当前基础分拣模型并行推进：短期先在统一评测下验证动作优化增益，随后把数字孪生用于需要更高定位精度的插管清洗任务。

\resultbox{第一年度已经完成从数据生产到现场运行的主要工程闭环，并观察到完整双臂长程任务能力。下一年度的核心不是推倒重来，而是把已有能力标准化量化、针对两个明确场景补强数据，并在同一协议下验证强化学习和高精度操作的增益。}
